\documentclass[11pt]{article}
\usepackage[a4paper,margin=25mm]{geometry}
\usepackage{amsmath,amssymb,amsthm,booktabs}
\usepackage[T1]{fontenc}
\usepackage{lmodern}
\usepackage[hidelinks]{hyperref}
\newtheorem{theorem}{Theorem}
\newtheorem{lemma}{Lemma}
\newcommand{\N}{\mathbb N}
\newcommand{\cdeg}{\mathsf c}
\title{Catenary degree is not determined by the minimum generator\\
\large A counterexample to Conjecture 00000002212}
\author{asfigh665}
\date{15 October 2026}
\begin{document}
\maketitle

\begin{abstract}
The minimally three-generated numerical semigroups
$H=\langle3,4,5\rangle$ and $K=\langle3,7,8\rangle$ have the same
minimum generator, but $\cdeg(H)=3$ and $\cdeg(K)=5$.
We prove the global upper bounds by terminating normalizations of every
factorization, and the matching lower bounds from explicit two-element fibers.
The entire argument, including the numerical-semigroup hypotheses and the
negation of the proposed functional dependence, is formalized in Lean 4.24.0.
\end{abstract}

\section{Definitions and the claim being refuted}
Write $\N=\{0,1,2,\ldots\}$. For positive generators $g=(g_1,g_2,g_3)$,
put
\[
 |x|=x_1+x_2+x_3,\qquad
 v_g(x)=g_1x_1+g_2x_2+g_3x_3,\qquad
 Z_g(n)=\{x\in\N^3:v_g(x)=n\}.
\]
For $x,y\in\N^3$, let $(x\wedge y)_i=\min(x_i,y_i)$ and define
\[
 d(x,y)=\max\{|x-x\wedge y|,\ |y-x\wedge y|\}.
\]
An $N$-chain in $Z_g(n)$ is a finite sequence in that fiber whose consecutive
terms have distance at most $N$. The catenary degree $\cdeg(\langle g\rangle)$
is the least $N$ such that every pair of factorizations of every element can be
joined by an $N$-chain. This is the standard catenary invariant: equivalently,
all distance-threshold graphs on factorization fibers are connected. The
problem's informal parenthetical ``diameter of the factorization graph'' is
used in this standard sense, not as ordinary unweighted graph diameter.

Conjecture 00000002212 asserts that for minimally three-generated numerical
semigroups this invariant is a function of the minimum generator. We disprove
even the existence of such a function, without imposing the extra demand
that it be explicit.

\begin{theorem}\label{thm:main}
The semigroups $H=\langle3,4,5\rangle$ and $K=\langle3,7,8\rangle$ are
minimally three-generated numerical semigroups, and
\[
\cdeg(H)=3,\qquad \cdeg(K)=5.
\]
Consequently, their catenary degrees cannot be determined solely by their
common minimum generator $3$.
\end{theorem}

\section{The numerical-semigroup hypotheses}
Both sets of nonnegative linear combinations contain $0$ and are closed
under addition. Every integer $n\ge3$ belongs to $H$: according to its
residue modulo $3$, it is a multiple of $3$, or $4$ plus a nonnegative
multiple of $3$, or $5$ plus a nonnegative multiple of $3$.
Similarly every $n\ge6$ belongs to $K$, using $0,7,8$ as the three residue
representatives. Thus both have finite complement in $\N$.

Each displayed generator has only its one-letter factorization. In $H$,
$3$ is smaller than $4$ and $5$, and neither $4$ nor $5$ is a multiple of $3$;
no sum involving a larger generator can equal a smaller one. Also $5-4=1$
is not in $H$. For $K$, neither $7$ nor $8$ is divisible by $3$, and
$8-7=1$ is not in $K$. These observations prove that all three displayed
generators are irreducible. Any other positive element is represented by
at least two displayed generators and hence is reducible. Each semigroup
therefore has exactly three minimal generators, as required.

\section{Global upper bounds}
The following moves preserve the represented integer. All are applied only
when the coordinates being reduced remain nonnegative.
\begin{center}
\begin{tabular}{@{}lll@{}}
\toprule
Semigroup & Replacement of factors & Distance \\
\midrule
$H$ & $4+4\longmapsto3+5$ & $2$ \\
    & $5+5\longmapsto3+3+4$ & $3$ \\
    & $4+5\longmapsto3+3+3$ & $3$ \\
\addlinespace
$K$ & $7+7\longmapsto3+3+8$ & $3$ \\
    & $8+8\longmapsto3+3+3+7$ & $4$ \\
    & $7+8\longmapsto3+3+3+3+3$ & $5$ \\
\bottomrule
\end{tabular}
\end{center}
The stated distances hold in every surrounding factorization: the common
coordinates cancel in the definition of $d$.

Let a factorization have coordinates $(a,b,c)$, where $b,c$ count the two
generators other than $3$. If $b\ge2$, use the first move; otherwise, if
$c\ge2$, use the second; otherwise, if $b=c=1$, use the third.
Each move strictly decreases $b+c$. Thus the procedure terminates for
\emph{every} factorization at a normal form with $b+c\le1$.

For $H$, the possible normal forms have values
\[
3a,\qquad 3a+4,\qquad 3a+5.
\]
They have distinct residues modulo $3$. The represented value determines
which form occurs and then determines $a$. Hence the normal form is unique
within each fiber. Any two factorizations of the same integer can be joined
by normalizing the first and reversing a normalization of the second.
Every step has distance at most $3$, proving $\cdeg(H)\le3$.

For $K$, the three normal-form values are $3a$, $3a+7$, and $3a+8$.
The same residue argument proves uniqueness. Its normalization moves have
distance at most $5$, so $\cdeg(K)\le5$.

\section{Sharp lower bounds and conclusion}
The fiber of $9$ for $H$ is exactly
\[
Z_{(3,4,5)}(9)=\{(3,0,0),(0,1,1)\}.
\]
Indeed, in $3a+4b+5c=9$, the case $c=0$ gives $(a,b)=(3,0)$, the case
$c=1$ gives $(a,b)=(0,1)$, and $c\ge2$ is impossible. The two
factorizations have distance $3$. Any chain between them must make a step
between the two distinct vertices, so no $N<3$ suffices.
Together with the upper bound, this gives $\cdeg(H)=3$.

Similarly,
\[
Z_{(3,7,8)}(15)=\{(5,0,0),(0,1,1)\}.
\]
For $c=0$, the equation $3a+7b=15$ has only $(a,b)=(5,0)$;
for $c=1$, it has only $(a,b)=(0,1)$; and $c\ge2$ is impossible.
These two factorizations have distance $5$, giving $\cdeg(K)\ge5$.
The global upper bound proves equality and completes Theorem~\ref{thm:main}.
If a function $f$ as conjectured existed, it would satisfy both $f(3)=3$
and $f(3)=5$, a contradiction.

\section{Formal verification and prior rejected work}
The file \texttt{Counterexample.lean} imports only \texttt{Std}.
The project is pinned to Lean 4.24.0. Its definitions
\texttt{distance}, \texttt{Chain}, \texttt{Bound}, and
\texttt{CatenaryDegree} implement precisely the definitions above.
\texttt{Admissible} requires positive distinct generators, unique one-letter
factorizations of all three generators, and an explicit conductor.
The submonoid properties are proved separately by \texttt{member\_zero}
and \texttt{member\_add}.

The unrestricted upper bounds use strong induction on $b+c$ in
\texttt{normalise} and \texttt{normalise\_h}. The sharp degrees are
\texttt{g\_degree\_three} and \texttt{h\_degree\_five}.
\texttt{counterexample} combines every hypothesis with the two exact degrees;
\texttt{conjecture\_false} formally negates the existence of the proposed
function. Standard kernel-checked arithmetic is used; there are no unproved
placeholders, additional axioms, or native decision procedures.
The supplied verification record lists the compiler result and final axiom
audit. Build with \texttt{lake build}, then audit with
\texttt{lake env lean Counterexample.lean}.

An earlier submission,
\href{https://github.com/The-Last-Math-Competition/The-Last-Math-Competition/pull/245\#issuecomment-5965086293}{PR \#245 (reviewer comment)},
was closed as invalid. The review identified redundant generators,
including $25=5\cdot5$ and $35=5\cdot7$, so its examples did not have
embedding dimension three. It also identified the absence of a Lean
definition of catenary degree, and the failure of finite stability windows
and the stated generic Betti argument to establish global upper bounds.
Those issues are addressed here by proved minimality, an explicit definition
of the invariant, and a terminating argument valid for all factorization
fibers. No finite search is used as evidence for a global upper bound.
\end{document}
